\chapter{The Research Program}
\label{ch:research-program}

The course has now built the mathematical foundation for a concrete research
program. The program is not merely to simulate clustered spiking networks until
interesting traces appear. It is to decide whether paired E/I clustered spiking
circuits realize a genuine mesoscale chaotic regime, to separate that regime
from finite-size multistability, and to explain why inhibitory clustering
changes the answer.

\section*{Learning Goals}

After this chapter you should be able to:
\begin{itemize}
    \item formulate phase diagrams over cluster strength, inhibitory clustering,
    heterogeneity, and network size;
    \item define the core observables needed to classify asynchronous,
    attractor, chaotic, and transient-chaotic regimes;
    \item compare the expected scaling behavior of E-only and paired E/I
    clustered networks;
    \item state what evidence would support mesoscopic chaos in a clustered
    spiking circuit.
\end{itemize}

\section{Central Questions}
\label{sec:ch10-central-questions}

The roadmap can be organized around four questions.
\begin{enumerate}
    \item Does mesoscopic chaos exist in clustered spiking circuits after
    finite-size effects are removed?
    \item Can one discriminate finite-size multistability from mesoscopic chaos
    using scaling, perturbation, autocorrelation, and Lyapunov diagnostics?
    \item Why should paired E/I clustered networks show stronger quenched
    mesoscale variability than E-only clustered networks with homogeneous
    inhibition?
    \item Can a mesoscale theory reproduce the same regimes and observables as
    the spiking network?
\end{enumerate}
These questions are linked. A positive answer to the first is not convincing
without the second. The third supplies the mechanism. The fourth makes the
claim mathematical rather than purely numerical.

\section{Parameter Axes and Phase Diagrams}
\label{sec:ch10-phase-diagrams}

The main phase diagrams should be built over cluster strength and inter-cluster
heterogeneity. A reduced two-parameter plane uses
\begin{equation}
    s \quad \text{or} \quad J_E^+
    \qquad \text{and} \qquad
    \kappa \quad \text{or} \quad g_{\mathrm{het}} .
    \label{eq:ch10-phase-axes}
\end{equation}
Here \(s\) is a cluster-strength parameter, often proportional to a potentiated
within-cluster weight such as \(J_E^+\). The parameter \(\kappa\), or
\(g_{\mathrm{het}}\) in older shorthand, controls quenched inter-cluster
heterogeneity, either in weights or in connectivity. Bare \(g\) is reserved for
inhibitory balance gains such as those in Chapter~\ref{ch:ei-clustered-networks}.
Additional axes include the inhibitory clustering ratio \(R\), the number of
clusters \(Q\), local population counts \(n_E,n_I\), and the finite-size
amplitude \(\xi=1/\sqrt{n_E}=(\gamma n_c)^{-1/2}\) at fixed local E/I fraction
\(n_c=n_E+n_I\), \(n_E=\gamma n_c\), \(n_I=(1-\gamma)n_c\). When total counts
are needed, write them explicitly as \(N_E^{\mathrm{tot}}=Qn_E\) and
\(N_I^{\mathrm{tot}}=Qn_I\).

\begin{table}[t]
\centering
\small
\begin{tabularx}{\textwidth}{@{}p{0.24\textwidth}p{0.32\textwidth}X@{}}
\toprule
\textbf{Phase label} & \textbf{Operational signature} & \textbf{Primary discriminator} \\
\midrule
Asynchronous irregular & Weak cluster modulation and no long cluster lifetimes & Low state occupancy structure and short autocorrelation. \\
Activated attractor & Stable active-cluster states with rare or absent switching & Switches disappear as \(n_E\) grows unless finite-size noise is retained. \\
Mesoscopic chaos & Persistent irregular cluster switching at small \(\xi\) & Positive mesoscale instability and switch-time plateau over the size ladder. \\
Transient chaos & Long irregular episodes followed by settling or restricted state exploration & Lifetime of the transient and state coverage must be measured separately. \\
\bottomrule
\end{tabularx}
\caption{Phase-diagram labels are assigned by diagnostics rather than by
example traces. Each grid point should carry the same observable bundle so that
neighboring regions can be compared reproducibly.}
\label{tab:ch10-phase-diagram}
\end{table}

\begin{table}[t]
\centering
\small
\begin{tabularx}{\textwidth}{@{}p{0.24\textwidth}p{0.28\textwidth}X@{}}
\toprule
\textbf{Quantity} & \textbf{Pedagogical default or sweep} & \textbf{Reason to record it} \\
\midrule
Architecture & E-only control and paired E/I clusters & Separates finite-size attractor switching from E/I quenched-disorder mechanisms. \\
Cluster count & \(Q\in\{10,20,40\}\), plus larger \(Q\) for regime C & Tests few-cluster sample-specific behavior versus many-unit cluster statistics. \\
Local population size & \(n_E:n_I=4:1\); sweep \(n_E\) by factors of two & Controls \(\xi=1/\sqrt{n_E}\) when testing finite-size scaling. \\
Connectivity construction & Fixed indegree or fully connected cluster-pair template from Chapter~\ref{ch:ei-clustered-networks} & Keeps mean input and realized indegrees from drifting across sweeps. \\
Cluster strength & Grid in \(J_E^+\) or \(s\); paired inhibitory strength via \(R\) & Locates asynchronous, attractor, chaotic, and transient regimes. \\
Quenched heterogeneity & Grid in \(\kappa\) or channel-specific \(\sigma_{\alpha\beta}^{\mathrm{het}}\) & Separates frozen inter-cluster disorder from finite-size sampling noise. \\
Integrator and recording & \(\Delta t=0.05\)--\(0.1\) ms; discard \(5\)--\(10\) s burn-in; record until at least dozens of transitions or report censoring & Prevents numerical and sampling choices from defining the phase label. \\
Rate extraction & \(\Delta_b=1\)--\(5\) ms bins; rate filter \(\tau_f=10\)--\(25\) ms; state detector with hysteresis and dwell time & Gives reproducible cluster rates, lifetimes, autocorrelations, and transition times. \\
Replicates & Multiple connectivity realizations and replayable stochastic streams per grid point & Separates within-realization dynamics from sample-to-sample quenched bias. \\
Perturbation ladder & Membrane-voltage perturbation, one-spike perturbation, whole-cluster perturbation, deterministic rate perturbation & Tests recovery, basin crossing, and infinitesimal sensitivity on comparable axes. \\
\bottomrule
\end{tabularx}
\caption{Canonical simulation template for the course project. The numbers are
pedagogical defaults for notebooks and smoke tests, not claims about the exact
settings used in any specific Tilo simulation. Publication-quality runs should
justify the grid, recording length, and thresholds from convergence checks.}
\label{tab:ch10-canonical-config}
\end{table}

A useful first classification has four regimes:
\begin{enumerate}
    \item \textbf{Asynchronous irregular regime.} Cluster strength is too weak
    to activate assemblies. Rates are low or weakly modulated, pairwise
    correlations are small, and no cluster has a long lifetime.
    \item \textbf{Attractor regime.} Cluster strength is large enough to create
    stable active-cluster states. At small \(\kappa\), switching is absent in
    the large-network limit or is driven by finite-size noise.
    \item \textbf{Chaotic regime.} Cluster strength and heterogeneity are large
    enough that the cluster-rate dynamics remains irregular as
    \(\xi\to 0\), with positive mesoscale Lyapunov exponent.
    \item \textbf{Transient-chaotic regime.} For very strong coupling, the
    network may show long chaotic transients before settling, or it may visit
    only a subset of active-cluster configurations. This regime should be
    treated as a hypothesis to test, not as an assumed phase.
\end{enumerate}

The roadmap makes two concrete trend predictions that should be checked on the
phase diagram rather than inferred from isolated traces. In paired E/I networks
with large cluster strength and large inter-cluster heterogeneity, the cluster
autocorrelation time \(\tau_c\) and cluster lifetime \(T_C\) should become
shorter as \(\kappa\) increases, as in random rate networks where stronger
disorder accelerates chaotic decorrelation. The same chaotic phase should be weakly
sensitive to network size when \(Q\) is fixed and \(n_c\) grows. In E-only
clustered networks, by contrast, the metastable window in cluster strength is
expected to be narrow, the timescales should be less controlled by \(\kappa\),
and transitions should be much more sensitive to cluster size because the
dominant escape source is finite-size noise.

\section{E-Only Versus E/I-Clustered Networks}
\label{sec:ch10-e-vs-ei}

The comparison between E-only and paired E/I networks is the core control
experiment. In an E-only clustered network, excitatory clusters are potentiated
but inhibition remains homogeneous. The expected mechanism for switching is
finite-size destabilization of stable attractors. As the number of neurons per
cluster grows at fixed number of clusters, the finite-size amplitude decreases
and transitions should become rarer or disappear.

In a paired E/I clustered network, each effective rate unit contains both an
excitatory and an inhibitory population. Quenched variability in inhibitory
inter-cluster structure can contribute strongly to input variability. A useful
schematic channel score is
\begin{equation}
    H_\alpha^{\mathrm{E/I}}
    =
    f_E
    (\sigma_{\alpha E}^{\mathrm{het}})^2
    j_{\alpha E}^2 r_E
    +
    f_I
    (\sigma_{\alpha I}^{\mathrm{het}})^2
    j_{\alpha I}^2 r_I.
    \label{eq:ch10-ei-heterogeneity-score}
\end{equation}
Here \(f_\beta=n_\beta/n_c\) is the local population fraction. This score ranks
channels after the connectivity and weight normalization have been fixed; it is
not by itself an implementable variance formula. The implementable quantity is
the empirical split of stored input currents into outside-cluster E and I
components, with within-realization temporal variability separated from
between-realization quenched bias.
Because inhibition-stabilized networks can have large \(|j_{\alpha I}|\), the
second term can dominate. The prediction is therefore qualitative as well as
quantitative: E/I clustering should preserve irregular cluster switching in
large clusters over a broader parameter range than E-only clustering.

\section{Observable Dashboard}
\label{sec:ch10-observable-dashboard}

Every simulation should produce a common set of observables. The point is to
avoid judging regimes from example traces alone.

\begin{table}[t]
\centering
\small
\begin{tabularx}{\textwidth}{@{}p{0.25\textwidth}p{0.28\textwidth}X@{}}
\toprule
\textbf{Observable family} & \textbf{Stored output} & \textbf{Question answered} \\
\midrule
Timescales & \(\tau_i\), \(\tau_c\), cluster lifetimes, switch rates & Are slow dynamics single-cell, cluster-level, or state-transition dominated? \\
State statistics & \(S_a(t)\), active-count distribution, transition table & Which active-cluster configurations are visited and how often? \\
Variability & Fano factors, current-component means and variances & Which source of input variability grows before transitions? \\
Dimensionality & Raw and residual participation ratios & Are transitions stereotyped paths or high-dimensional excursions? \\
Instability & Common-noise perturbation growth and mesoscale Lyapunov estimates & Do nearby trajectories recover or separate at small \(\xi\)? \\
\bottomrule
\end{tabularx}
\caption{Observable dashboard for each phase-diagram grid point. The same
bundle should be produced for spiking simulations, deterministic mesoscopic
runs, and synthetic validation traces.}
\label{tab:ch10-observable-dashboard}
\end{table}

\paragraph{Single-neuron autocorrelation time \(\tau_i\).}
The single-neuron observable starts from binned spike counts or instantaneous
single-neuron rates \(r_i^{(u)}(t)\) on trial \(u\), often using 5 ms bins. Its
raw trial-averaged autocorrelation is corrected by a shift predictor. For
positive lags \(0\le \tau<T\), all finite-window averages use the valid overlap
\([0,T-\tau]\), not the unavailable part of the trace. The shift-corrected
autocorrelation is
\begin{equation}
    a_i^{\mathrm{corr}}(\tau)
    =
    \frac{1}{n_{\mathrm{tr}}}
    \sum_{u=1}^{n_{\mathrm{tr}}}
    \frac{1}{T-\tau}
    \int_0^{T-\tau}
    r_i^{(u)}(t)r_i^{(u)}(t+\tau)\,dt
    -
    a_i^{(0)}(\tau),
    \label{eq:ch10-single-neuron-autocorrelation}
\end{equation}
where
\begin{equation}
    a_i^{(0)}(\tau)
    =
    \frac{1}{n_{\mathrm{tr}}}
    \sum_{u=1}^{n_{\mathrm{tr}}}
    \frac{1}{T-\tau}
    \int_0^{T-\tau}
    r_i^{(u)}(t)r_i^{(\pi(u))}(t+\tau)\,dt
    \label{eq:ch10-shift-predictor}
\end{equation}
uses a random permutation \(\pi\) of trial labels. The central spike-count peak
at \(\tau=0\) is the Poisson-like same-bin contribution of discrete spikes; it
must be removed before measuring the slow rate autocorrelation. Let
\(a_i^{\mathrm{rate}}(\tau)\) denote this central-peak-removed function, and
measure HWHM on the normalized autocorrelation
\[
    \rho_i^{\mathrm{rate}}(\tau)
    =
    \frac{a_i^{\mathrm{rate}}(\tau)}
    {a_i^{\mathrm{rate}}(0^+)}.
\]
The single-neuron time \(\tau_i\) is
\begin{equation}
    \tau_i
    =
    \inf\{\tau>0:\rho_i^{\mathrm{rate}}(\tau)\le 1/2\}.
    \label{eq:ch10-single-neuron-hwhm}
\end{equation}
This keeps a fast counting artifact from being mistaken for a slow metastable
or chaotic timescale.

\paragraph{Cluster autocorrelation time \(\tau_c\).}
For a cluster rate \(r_c(t)\), define
\begin{equation}
\begin{aligned}
    C_c(\tau)
    &=
    \frac{1}{T-\tau}
    \int_0^{T-\tau}
    \delta r_c(t)\delta r_c(t+\tau)\,dt,
    \\
    \delta r_c(t)
    &=
    r_c(t)-\overline r_c,
    \qquad
    \overline r_c
    =
    \frac{1}{T}\int_0^T r_c(t)\,dt .
\end{aligned}
    \label{eq:ch10-autocorrelation}
\end{equation}
The cluster HWHM is also measured on the normalized autocorrelation
\(\rho_c(\tau)=C_c(\tau)/C_c(0)\). The cluster autocorrelation time
\(\tau_c\) is the smallest positive lag at which this normalized
autocorrelation has fallen to one half:
\begin{equation}
    \tau_c
    =
    \inf\{\tau>0:\rho_c(\tau)\le 1/2\}.
    \label{eq:ch10-cluster-hwhm}
\end{equation}
This is the HWHM autocorrelation time used to compare clustered networks with
SCS and self-coupled rate models. Unlike \(\tau_i\), it is computed directly
from the population rate of a cluster and does not contain the single-spike
central peak.

\paragraph{Cluster lifetimes.}
Choose a threshold \(r_0\) separating inactive and active cluster rates. A
cluster lifetime \(T_C\) is the length of a contiguous interval on which
\begin{equation}
    r_c(t)>r_0.
    \label{eq:ch10-cluster-lifetime}
\end{equation}
The lifetime distribution distinguishes rare noise-driven escapes from faster
chaotic switching, especially when measured across network sizes.

\paragraph{Firing-rate distributions and active-cluster counts.}
The marginal distribution of \(r_c(t)\) reports whether clusters are mostly
silent, bimodal, broadly distributed, or saturated. The active-cluster count
\begin{equation}
    M(t)=\sum_{c=1}^{Q}\mathbf{1}\{r_c(t)>r_0\}
    \label{eq:ch10-active-cluster-count}
\end{equation}
summarizes how many assemblies participate at a time. Attractor regimes often
have preferred counts; chaotic regimes may have broader or parameter-dependent
count distributions.

\paragraph{Fano factors.}
For a spike count \(N_i(T)\) in a window of length \(T\), the Fano factor is
\begin{equation}
    F_i(T)
    =
    \frac{\operatorname{Var}[N_i(T)]}{\mathbb{E}[N_i(T)]}.
    \label{eq:ch10-fano}
\end{equation}
A Poisson process has \(F=1\). Deviations reflect rate modulation, correlations,
and state switching. Fano factors are especially useful because metastable
activity was originally linked to changes in variability during cortical state
transitions and task engagement.

\paragraph{Dimensionality.}
Let \(C\) be the pairwise covariance or correlation matrix of cluster firing
rates. The participation-ratio dimensionality is
\begin{equation}
    d
    =
    \frac{\operatorname{Tr}(C)^2}{\operatorname{Tr}(C^2)}.
    \label{eq:ch10-dimensionality}
\end{equation}
If activity lies along one dominant mode, \(d\) is small. If many cluster modes
participate, \(d\) is larger. A chaotic regime should be characterized not only
by positive Lyapunov exponent but also by the dimension of the activity it
explores.

\paragraph{Transition-triggered variability and cross-correlations.}
Detect isolated off-to-on and on-to-off transitions of a focus cluster \(c\):
within \([t_m-T_C,t_m+T_C]\), no other cluster should cross the activity
threshold. For transition times \(t_m\), compute
\begin{equation}
    \operatorname{TTA}_{\sigma^2}(\tau)
    =
    \frac{1}{M}\sum_{m=1}^M \sigma^2(t_m+\tau),
    \label{eq:ch10-tta-variability}
\end{equation}
where \(\sigma^2\) should be split into spatial input-current variance
components across the excitatory neurons of cluster \(c\):
\begin{equation}
    \sigma_{\mathrm{in}}^E,\quad
    \sigma_{\mathrm{in}}^I,\quad
    \sigma_{\mathrm{out}}^E,\quad
    \sigma_{\mathrm{out}}^I .
    \label{eq:ch10-transition-sigma-components}
\end{equation}
For \(x\in\{\mathrm{in},\mathrm{out}\}\) and \(\beta\in\{E,I\}\), each
component means
\begin{equation}
    \left(\sigma_x^\beta(t)\right)^2
    =
    \frac{1}{n_E}
    \sum_{i=1}^{n_E}
    \left(
        I_{x,i}^{\beta}(t)-\overline I_x^\beta(t)
    \right)^2 .
    \label{eq:ch10-transition-sigma-definition}
\end{equation}
The matching mean current components should also be triggered. This asks which
source drives the transition: within-cluster excitation, within-cluster
inhibition, outside-cluster excitation, or outside-cluster inhibition.

For a specific state transition \(A\to B\), also project the triggered rate
trajectory onto the transition direction
\begin{equation}
    e_{AB}=\frac{r_B-r_A}{\|r_B-r_A\|},
    \qquad
    z_{\parallel}(t)=e_{AB}^{\mathsf T}(r(t)-r_A),
    \label{eq:ch10-transition-direction}
\end{equation}
and compare it with the orthogonal displacement
\begin{equation}
    z_{\perp}^2(t)
    =
    \|r(t)-r_A\|^2-z_{\parallel}^2(t).
    \label{eq:ch10-transition-orthogonal}
\end{equation}
The predicted signature of a directional transition is a pre-crossing increase
along \(r_B-r_A\) without a comparable increase in orthogonal directions.
Cross-correlations between rate changes and the four variance components test
whether fluctuations precede transitions or merely follow them.

\paragraph{Mesoscale Lyapunov exponents.}
Estimate \(\lambda_{\max}\) from the cluster-rate dynamics using perturbation
growth or tangent dynamics:
\begin{equation}
    \lambda_{\max}
    =
    \lim_{t\to\infty}
    \frac{1}{t}
    \log
    \frac{\|\delta r(t)\|}{\|\delta r(0)\|}.
    \label{eq:ch10-lyapunov}
\end{equation}
This observable is the most direct test for rate chaos. It should be reported
with the network-size scaling of the same parameter point.

\section{Scaling Experiments}
\label{sec:ch10-scaling-experiments}

The roadmap proposes three complementary scalings.

\paragraph{Fixed cluster size, increasing number of clusters.}
Hold \(n_c\) fixed and increase \(Q\). Finite-size effects remain because each
cluster is still estimated from the same number of neurons. This scaling is
useful for studying many-cluster statistics but cannot eliminate finite-size
switching.

\paragraph{Fixed number of clusters, increasing cluster size.}
Hold \(Q\) fixed and increase \(n_c\). With fixed local E/I fraction
\(n_E=\gamma n_c\), this scaling suppresses
\(\xi=1/\sqrt{n_E}=(\gamma n_c)^{-1/2}\). It is the decisive test for whether
transitions in E-only networks are finite-size driven. In paired E/I networks,
persistence of switching under this scaling supports a quenched mesoscale
mechanism.

\paragraph{Increasing both clusters and cluster size.}
Let both \(Q\) and \(n_c\) grow, for example
\(Q=(N^{\mathrm{tot}})^\rho\) and
\(n_c=(N^{\mathrm{tot}})^{1-\rho}\) with \(0<\rho<1\). This is the scaling that
connects the spiking network to a many-unit dynamic mean-field theory. The
finite-size amplitude vanishes while the number of quenched cluster couplings
grows.

\section{Simulation and Theory Workflow}
\label{sec:ch10-workflow}

A rigorous project loop should have the following structure.
\begin{enumerate}
    \item Choose architecture: E-only or paired E/I clusters; fully connected,
    neuron-level Erdos-Renyi, or cluster-level Erdos-Renyi inter-cluster
    variability.
    \item Choose parameters: \(J_E^+\), \(R\), depression factors
    \(J_E^-,J_I^-\), heterogeneity amplitude \(\kappa\) or \(g_{\mathrm{het}}\),
    number of clusters \(Q\), and neurons per cluster \(n_c\).
    \item Simulate spiking dynamics and record spike trains, filtered rates,
    input-current components, and cluster labels.
    \item Compute the observable dashboard: autocorrelations, HWHM times,
    lifetimes, rate distributions, active-cluster counts, Fano factors,
    dimensionality, transition-triggered variability, cross-correlations, and
    perturbation growth.
    \item Repeat across network sizes to separate \(\xi\)-dependent behavior
    from \(\kappa\)-dependent behavior.
    \item Compare the spiking network to the mesoscale theory and, where
    possible, to a dynamic mean-field self-consistency calculation.
\end{enumerate}

\paragraph{Minimum run products.}
Every run should write a machine-readable manifest and a compact analysis
bundle. The manifest records architecture, \(Q\), \(n_E\), \(n_I\), all weight
and depression factors, fixed indegrees, heterogeneity amplitudes, random
seeds, integration step, burn-in, recording length, filter widths, threshold
rules, and software version. The analysis bundle stores cluster-rate traces,
state vectors \(S_a(t)\), transition tables, autocorrelations, HWHM estimates,
Fano curves, dimensionality curves, input-current variance decompositions, and
perturbation-growth summaries. A phase diagram assembled without these run
products is not reproducible enough to support a mechanism claim.

The research loop is:
\begin{enumerate}
    \item choose the model architecture, \(Q,n_E,n_I\), cluster strength,
    inhibitory clustering, and heterogeneity grid;
    \item run simulations with manifests and replayable stochastic streams;
    \item build the common dashboard of rates, states, autocorrelations,
    transition statistics, perturbation growth, and current decompositions;
    \item repeat the dashboard across the \(n_E\), \(Q\), and \(\kappa\)
    scaling experiments;
    \item compare the resulting phase labels with deterministic mesoscopic or
    DMFT predictions before refining the next sweep.
\end{enumerate}

The workflow is intentionally redundant. If a parameter point is truly chaotic,
it should not rely on one fragile signature. It should show persistent switching
as \(\xi\to 0\), positive mesoscale Lyapunov exponent, smooth broad
autocorrelation, and a stable location in the phase diagram.

\section{Contribution-Level Research Questions}
\label{sec:ch10-contribution-questions}

The program naturally separates into publishable contributions.

\paragraph{Mechanism.}
Can the channel score in \eqref{eq:ch10-ei-heterogeneity-score}, together with
empirical current decompositions, predict where E/I clustering amplifies
quenched disorder enough to destabilize attractors?

\paragraph{Diagnostics.}
Can switch-time plateaus, perturbation divergence, zero-lag autocorrelation
shape, and mesoscale Lyapunov exponents be made consistent on the same
parameter grid?

\paragraph{Theory.}
Can one derive a dynamic mean-field theory for paired E/I cluster units in the
large-\(Q\), large-\(n_c\) limit and obtain a self-consistency equation for
\(C(\tau)\)?

\paragraph{Comparison.}
Can the E-only and E/I-clustered networks be matched in mean firing rate and
cluster lifetime while differing in their large-\(n_c\) scaling? This would be
a strong demonstration that similar phenomenology can arise from different
mechanisms.

\paragraph{Neuroscience.}
Can the model explain why metastable cortical activity is useful for state-
dependent variability, expectation, attention, and decision timing while also
predicting when such activity should recover from perturbation versus diverge?

\section{Checkpoint Exercises and Notebook Prompts}
\label{sec:ch10-checkpoint-exercises}

The fastest route into the project is a sequence of notebooks that each
produces one reusable diagnostic artifact. The checkpoints below are stated as
implementation exercises because the research bottleneck is usually not the
definition of an observable; it is making that observable robust enough to
survive parameter sweeps.

\paragraph{Checkpoint 1: spike-to-rate extraction.}
Implement the binner and exponential filter in
\eqref{eq:ch09-discrete-rate-filter}. Validate units with a homogeneous
Poisson process: the mean filtered rate should equal the imposed rate, and the
Fano factor should be near one before any state modulation is introduced. The
notebook output is a rate table indexed by time, cluster, and cell type.

\paragraph{Checkpoint 2: state detector and lifetime table.}
Use synthetic two-state traces with known dwell times to tune threshold,
hysteresis, and minimum-dwell rules. Then apply the detector to one E-only and
one paired E/I simulation at matched mean rate. The notebook output is a table
of on/off events, active-state vectors, active-cluster counts, and lifetime
distributions with threshold-sensitivity bands.

\paragraph{Checkpoint 3: autocorrelation, HWHM, and kink estimator.}
Compute valid-overlap autocorrelations, HWHM times, and the small-lag kink fit
from \eqref{eq:ch09-kink-fit}. Validate on three signals: Poisson counts,
Ornstein-Uhlenbeck rates, and a smooth deterministic chaotic or quasiperiodic
test trace. The notebook output is a plot of \(k_\xi\) and
\(\tau_{\mathrm{HWHM}}\) versus local population size.

\paragraph{Checkpoint 4: transition-triggered variability.}
For isolated transitions, compute triggered means and variances of
\(I_{\mathrm{in}}^E\), \(I_{\mathrm{in}}^I\), \(I_{\mathrm{out}}^E\), and
\(I_{\mathrm{out}}^I\), plus transition-triggered Fano factors with
rate-matched control windows. The notebook output is an aligned transition
panel that states whether the leading signature is increased excitation,
decreased inhibition, increased spatial variance, or no reliable precursor.

\paragraph{Checkpoint 5: dimensionality of state exploration.}
Build the cluster-rate covariance matrix for whole recordings and for
transition-aligned windows. Report participation-ratio dimensionality before
and after subtracting the triggered mean path. The notebook output is a
dimension curve over \(J_E^+\), \(R\), and \(\kappa\), with bootstrap intervals
over network realizations.

\paragraph{Checkpoint 6: finite-size scaling decision.}
Run a size ladder at fixed \(Q\) and fixed cluster-level parameters. Fit both
templates in \eqref{eq:ch09-size-scaling-fits}, report right-censored lifetimes
explicitly, and show whether the same grid point keeps its autocorrelation
timescale and active-state statistics as \(n_E\) grows. The notebook output is
a finite-size decision plot: escape-like divergence, plateau-like persistence,
or inconclusive mixed behavior.

\paragraph{Checkpoint 7: quenched-disorder sweep.}
At large \(n_E\), sweep \(\kappa\) while saving each connectivity manifest.
Decompose variability into within-realization temporal variance and
between-realization quenched bias using
\eqref{eq:ch09-time-realization-variance-split}. The notebook output is a
phase strip that shows where increasing \(\kappa\) creates temporal switching
rather than merely selecting static high-rate clusters.

\paragraph{Checkpoint 8: Lyapunov implementation.}
Validate the Benettin or QR code on a stable linear system and on a reference
chaotic rate model. Then apply it to the deterministic mesoscopic cluster
model. In spiking simulations, run the common-noise perturbation ladder and
compare projected growth rates across \(n_E\). The notebook output is a
convergence panel for \(\widehat\lambda_{\max}(T)\), including sensitivity to
time step, renormalization interval, perturbation size, and norm.

\paragraph{Checkpoint 9: spiking-to-mesoscopic comparison.}
Choose a small set of matched parameter points and compare spiking observables
with the deterministic or stochastic mesoscopic reduction: mean rates,
autocorrelation widths, active-count distributions, transition rates, and
Lyapunov estimates. The notebook output is a mismatch table that identifies
which terms in the mesoscopic theory need refinement before the model can be
used predictively.

\begin{conceptquestions}
    \item Which checkpoint would fail first if finite-size hopping were being
    mistaken for mesoscopic chaos, and why?
    \item Which observable is most sensitive to the threshold \(r_0\), and how
    would you design a robustness plot that does not hide that sensitivity?
    \item In a paired E/I network with persistent switching at large \(n_E\),
    what extra evidence is needed before calling the regime deterministic
    chaos?
\end{conceptquestions}

\section{What Would Count as Evidence?}
\label{sec:ch10-evidence}

The strongest evidence for mesoscopic chaos in paired E/I clustered spiking
networks would be the following package:
\begin{enumerate}
    \item Switching persists as \(n_c\) increases and
    \(\xi=(\gamma n_c)^{-1/2}\to 0\).
    \item The switch-time curve approaches a plateau controlled by \(\kappa\).
    \item Small perturbations grow with a positive mesoscale Lyapunov exponent.
    \item The cluster-rate autocorrelation keeps a broad timescale without a
    finite-size zero-lag kink.
    \item The same parameter regime appears in a mesoscale or DMFT reduction.
    \item The effect is stronger in paired E/I clusters than in E-only clusters
    with homogeneous inhibition.
\end{enumerate}
This package would support the claim that paired E/I spiking networks implement
a biologically plausible version of SCS-like rate chaos at the mesoscale.

\section{Chapter Summary}
\label{sec:ch10-summary}

The research program maps clustered spiking networks by cluster strength,
quenched heterogeneity, inhibitory clustering ratio, and network size. Its
observables include single-neuron \(\tau_i\), cluster HWHM autocorrelation
time \(\tau_c\), cluster lifetimes, firing-rate distributions,
active-cluster counts, Fano factors, dimensionality, transition-triggered
variability, scaling properties, and mesoscale Lyapunov exponents. The key
comparison is E-only versus paired E/I clustering. The key classification is
asynchronous, attractor, chaotic, or transient-chaotic. The key claim is
mechanistic: paired inhibitory structure can amplify quenched inter-cluster
variability enough to produce macroscopic switching without finite-size noise.
The canonical simulation template and checkpoint notebooks turn that claim into
a rotation plan: build the rate and state pipeline, validate every diagnostic
on synthetic controls, sweep \(\xi\) and \(\kappa\), and require Lyapunov,
autocorrelation, transition, dimensionality, and scaling evidence to agree
before assigning a chaotic phase label.
